\Chapter{Summary}

The aim of this thesis was to analyze video game data obtained from the Steam \cite{steam} platform and to formulate different classification problems using machine learning methods.
During the research a real, heterogeneous dataset was processed, which contains numerical variables, categorical metadata and textual descriptions.
One of the main objectives of the work was to demonstrate how these different types of information can be used in predictive modelling tasks.

The first step of the processing was the cleaning and standardization of the data.
This included the handling of missing values, the unification of variable types,
and the construction of features suitable for the models.
Transforming categorical fields into binary representations,
as well as the TF--IDF based vectorization of textual descriptions,
made it possible for the different types of data to appear in a common feature space.

The thesis examined two different classification problems.
The first task was the estimation of the success of video games
based on a label derived from user reviews.
In this case structured metadata and numerical variables formed the input of the predictive models.
The analysis showed that certain functional characteristics and platform services of games
may be related to the success reflected by user feedback.

The second examined task was the automatic separation of multiplayer and singleplayer games.
In this case textual descriptions and structured tags proved to be the most informative features.
The textual representation made it possible for the models
to distinguish between the two categories
based on expressions referring to game mechanics and functions.

The comparison of the different classification models showed
that both linear and nonlinear methods are capable of producing interpretable and stable predictions.
The advantage of linear models primarily appeared in interpretability
and in the efficient handling of high-dimensional feature spaces,
while ensemble methods in some cases showed better predictive performance.
Overall, the results indicate that the metadata
and textual information related to video games
are suitable for being used as inputs of machine learning models.

One important conclusion of the thesis is
that the combination of structured and textual features
can significantly increase the information content of classification models.
Game descriptions and metadata do not only contain marketing information,
but also implicit signals related to game mechanics and user experience,
which may also be useful for predictive models.

The presented approach can be further developed in several directions.
One possible direction of future research is the application
of more advanced natural language processing methods,
for example the use of word embeddings or neural language models.
These may be capable of capturing the deeper semantic structure of texts,
which could further improve classification performance.

Another possible development direction is the application
of more complex machine learning models,
for example deep neural networks or more advanced ensemble methods.
In addition, the concept of success could be modelled in a more detailed way,
for example by using multiclass or regression approaches,
where the target variable is not binary
but continuous or consists of several categories.

Overall, the thesis demonstrated
that the video game data available on the Steam platform
are suitable for data-driven analysis and predictive modelling tasks.
The applied methods proved to be useful not only from the perspective of predictive performance,
but also made it possible to better understand the relationships
between the characteristics of video games
and user feedback.